\section{Resultados esperados}

\begin{table}[H]
	\centering
	\label{tabla:resultadosEsperados}
	\caption{Relación entre los objetivos específicos y los resultados esperados}
	\begin{tabular}{|p{5.5cm}|p{7.5cm}|}
		\hline
			\cellcolor[gray]{0.9} \textbf{Objetivo específico} & \cellcolor[gray]{0.9}\textbf{Resultado esperado} \\
		\hline
		
1. Analizar los requerimientos reales de una PYME del sector alimentos.

 & 
Un documento de requisitos con historias de usuario, criterios de aceptación y las conclusiones de las entrevistas con la panificadora caso de estudio.

\\
		\hline
		
2. Diseñar la arquitectura del sistema y la experiencia de usuario.

 & 
Un documento técnico con el modelo Entidad-Relación, diagramas UML de los procesos clave y los 18 mockups (prototipos visuales) de la interfaz de usuario.

\\
		\hline
	    
3. Desarrollar los módulos del sistema, incluyendo el asistente de consultas inteligentes.

& 
El código fuente completo del sistema, versionado en Git, listo para funcionar y con el módulo de consultas en lenguaje natural integrado.

        \\
        \hline
		
4. Validar que el sistema funciona, es usable y es accesible.

 & 
Un informe de pruebas con los resultados de las pruebas funcionales, de usabilidad (heurísticas de Nielsen) y de accesibilidad (WCAG AA).

\\
        
        \hline
	\end{tabular}
\end{table}
